\section{Immutable Parent Pointers}
\label{sec:immutable-parent-pointers}

The immutable parent pointer is the foundational structural primitive of the Recursive Spawning Protocol. It is the single mechanism that transforms the delegation chain from a forensic reconstruction attempted after the fact into a runtime-verifiable artifact — one that is immutable by architectural constraint, auditable at any runtime point, and structurally enforcement through the Fact Layer write protocol. It renders all modes of autonomous drift structurally impossible regardless of operational urgency, architectural complexity, or adversarial sophistication of any machine actor in the chain.

This section specifies the parent pointer formally (\S\ref{sec:parent-pointer-formal}), defines the chain operator and establishes the base case for the root human anchor (\S\ref{sec:chain-definition}), specifies immutability enforcement through the Fact Layer write protocol (\S\ref{sec:fact-layer-enforcement}), describes the Purpose Layer's spawn authorization role (\S\ref{sec:purpose-layer-authorization}), and provides the complete chain traversal and verification algorithms with correctness arguments (\S\ref{sec:chain-traversal}).

% ------------------------------------------------------------------------- %
\subsection{Formal Definition: ParentPointer}
\label{sec:parent-pointer-formal}

\begin{dfn}[Parent Pointer]
\label{dfn:parent-pointer}
Let $\mathcal{N}$ be the set of all agent nodes in a \bxthree{} system implementing the Recursive Spawning Protocol, and let $P$ be the set of all Purpose Layer actors. The \emph{parent pointer} $\alpha(n)$ is defined for every node $n \in \mathcal{N}$ as the total function

\begin{equation}
\alpha : \mathcal{N} \rightarrow P \cup \{\bot\}
\label{eq:alpha-function}
\end{equation}

where $\alpha(n) = p$ if $p \in P$ is the Purpose Layer actor that authorized the spawn event resulting in $n$'s activation, and $\alpha(n) = \bot$ if and only if $n$ is the root human anchor established at system initialization. The pointer is established exactly once at node provisioning via a Fact Layer atomic write operation and satisfies the following immutability property:

\begin{equation}
\forall n, p : \bigl(\alpha(n) = p \ \land \ \textit{provisioned}(n)\bigr) \implies \neg\exists\ \textit{modify}(\alpha(n))
\label{eq:parent-pointer-immutable}
\end{equation}

No actor — including the Purpose Layer after provisioning, the Bounds Engine, any administrative process, or any external principal — may execute an operation that overwrites, redirects, or deletes $\alpha(n)$ after the provisioning write is confirmed. The Fact Layer write protocol rejects any such attempt at the protocol level, before execution reaches the ledger.
\end{dfn}

The parent pointer is not a reference variable held by the parent node toward its child. It is a Fact Layer property intrinsic to the child node, established at provisioning time and stored in the node's sealed memory envelope. The directionality is architecturally critical: the child carries $\alpha(n)$ as an immutable attribute of itself, not as a mutable external reference held by the parent. This means the chain is traversable from any node to its complete lineage regardless of the operational state of any other node in the chain. A parent that has terminated, crashed, or been decommissioned does not affect the child's ability to reference its parent through $\alpha(n)$.

The parent pointer is distinct from a conventional access-controlled pointer in a second critical respect. Access control enforces immutability by restricting which principals may issue modification operations; if sufficient privilege is acquired through credential theft, software vulnerability, or administrative compromise, the access control barrier is circumvented. The Fact Layer write protocol enforces immutability by making the modification operation structurally impossible — there is no defined mechanism for processing a $\textit{modify}(\alpha(n))$ request regardless of any privilege held. There is no privileged override, no emergency bypass, and no administrative role with the capability to alter $\alpha(n)$ after provisioning. This is not a stronger access control policy; it is a different architectural category.

\begin{dfn}[Spawn Event]
\label{dfn:spawn-event}
A \emph{spawn event} is the 5-tuple

\begin{equation}
\mathrm{spawn}(c,\ p,\ \sigma_c,\ t,\ \phi)
\label{eq:spawn-event}
\end{equation}

where $c$ is the newly created child node, $p \in P$ is the authorizing Purpose Layer actor, $\sigma_c$ is the authorized execution scope, $t$ is the wall-clock timestamp, and $\phi$ is the cryptographic signature produced by $p$ over the event payload using $p$'s Purpose Layer private key. The signature $\phi$ attests that $p$ has authorized the spawn and binds $c$ to $p$ in the forensic ledger. The Fact Layer verifies $\phi$ before accepting the spawn event for ledger recording.
\end{dfn}

The spawn event is the sole legitimate mechanism for assigning a parent pointer. There is no other path by which $\alpha(c)$ can be set. This eliminates any ambiguity about the origin of a node's parent pointer: if $\alpha(c) = p$ appears in the ledger, it exists because a spawn event bearing $p$'s cryptographic signature was recorded, and no other origin is possible.

% ------------------------------------------------------------------------- %
\subsection{Chain Definition and Root Human Base Case}
\label{sec:chain-definition}

The chain operator $\Chain{\cdot}$ constructs the complete delegation lineage from any node to the human anchor by recursively applying $\alpha$. It is the primary tool for chain verification, audit, and accountability attribution.

\begin{dfn}[Chain Operator]
\label{dfn:chain-operator}
For any node $a \in \mathcal{N}$, the \emph{chain} $\Chain{a}$ is defined recursively as:

\begin{equation}
\Chain{a} =
\begin{cases}
\{\, \ParentPointer(p,\ a) \, \} \cup \Chain{p} & \text{if } \alpha(a) = p \neq \bot \\[10pt]
\emptyset & \text{if } \alpha(a) = \bot \quad \text{(root human anchor)}
\end{cases}
\label{eq:chain-operator}
\end{equation}

where $\ParentPointer(p,\ a)$ denotes the ordered pair $(p,\ a)$ establishing that $p$ is the immutable parent of $a$. The empty set $\emptyset$ returned in the base case indicates that the root human has no parent pointer — it is the architectural terminus of all chains.

Equivalently, expanding the recursion:

\begin{equation}
\Chain{a} = \{\, a,\ \alpha(a),\ \alpha^{(2)}(a),\ \alpha^{(3)}(a),\ \ldots,\ n_0 \,\}
\label{eq:chain-expanded}
\end{equation}

where $n_0$ is the unique root satisfying $\alpha(n_0) = \bot$.
\end{dfn}

The base case is structurally non-negotiable: the root human anchor is the only node in any RSP chain with $\alpha(n) = \bot$. No machine actor can serve as a chain root — this is enforced by the Fact Layer pre-commit hook, which rejects any spawn event where the parent node does not satisfy the human anchor reachability predicate.

\begin{dfn}[Root Human Anchor]
\label{dfn:root-human}
The root human anchor $h \in P$ is the unique Purpose Layer actor designated at system initialization as the terminus of all delegation chains. It satisfies $\alpha(h) = \bot$ and $\mathrm{depth}(h) = 0$. The root human is non-collapsing and non-substitutable: as chains grow in depth through successive spawn events, the root remains fixed. When a child node $n_i$ spawns a grandchild $n_{i+1}$, the root for $n_{i+1}$ is the same root as for $n_i$; the root does not advance to intermediate machine actors.
\end{dfn}

The depth of any node $n$ is defined recursively:

\begin{equation}
\mathrm{depth}(n) \equiv
\begin{cases}
0 & \text{if } \alpha(n) = \bot \\
1 + \mathrm{depth}(\alpha(n)) & \text{otherwise}
\end{cases}
\label{eq:depth-definition}
\end{equation}

\medskip
\noindent\textbf{Properties of the chain operator.}

\begin{itemize}
    \item[(C1) \textbf{Determinism}] For any node $a$, $\Chain{a}$ is uniquely determined. There is no alternative chain, no ambiguous parentage, and no non-deterministic branching. $\alpha$ is a total function with unique values, so the chain is uniquely determined at provisioning and remains so at all runtime points.
    \item[(C2) \textbf{Stability}] The immutability of $\alpha$ guarantees that $\Chain{a}$ at any runtime point $t$ is identical to $\Chain{a}$ at $a$'s provisioning moment. The chain is not reconstructed from mutable logs — it is directly computable by iterating $\alpha$ from any node until reaching $\bot$.
    \item[(C3) \textbf{Compositionality}] For any two nodes $a$ and $b$ in the same RSP chain, either $\Chain{a} \subseteq \Chain{b}$ or $\Chain{b} \subseteq \Chain{a}$. The chains are nested rather than intersecting. No cross-links or merging paths exist in a valid RSP chain, which follows directly from the uniqueness constraint on parent pointers.
    \item[(C4) \textbf{Acyclicity}] $\Chain{a}$ is cycle-free by construction. Each recursive step applies $\alpha$ to the parent node, which is strictly closer to the root by depth. The Fact Layer rejects any spawn event that would introduce a cycle, and $\alpha$ is total. Every node has exactly one parent unless it is the root human.
\end{itemize}

% ------------------------------------------------------------------------- %
\subsection{Immutability Enforcement: Fact Layer Write Protocol}
\label{sec:fact-layer-enforcement}

The immutability of $\alpha(n)$ is enforced by the Fact Layer write protocol — the \bxthree{} mechanism that makes certain state transitions structurally irreversible. The write protocol defines which operations are valid for each node property at each node state; it is not access control with additional layers, but a protocol-level constraint on the state space itself.

\medskip
\noindent\textbf{Write protocol definition.} For the parent pointer $\alpha(n)$, the write protocol defines exactly one valid write operation:

\begin{enumerate}
    \item[$\textbf{W1}$] $\mathrm{WriteParentPointer}(c,\ p)$: The initial provisioning write, executed exactly once during a spawn event, recording $\alpha(c) = p$.
    \item[$\textbf{W2}$] No subsequent write operation targeting $\alpha(n)$ for any $n$ where $\textit{provisioned}(n) = \text{true}$.
\end{enumerate}

The Fact Layer write protocol processes write requests as follows: (1) authenticate the requestor's identity; (2) verify the request type against the defined write operation set; (3) execute the write if the type is W1 and all preconditions are satisfied; or (4) reject with error code \texttt{E-NOT-PERMITTED} if the type is W2 or any other undefined operation targeting $\alpha(n)$. Steps (1) and (2) are applied to every write request without exception. There is no bypass path, no debug mode, and no emergency override that disables step (2).

The rejection is unconditional and source-equivalent:

\begin{itemize}
    \item \textbf{Failure resilience.} If the system is in a degraded operational state — partial network partition, Bounds Engine timeout, Purpose Layer unavailability — the Fact Layer write protocol for $\alpha(n)$ immutability remains in force. The modification request is rejected regardless of system state, because the write protocol is a local constraint at the Fact Layer, not a distributed agreement protocol that could be blocked by network failure.
    \item \textbf{Administrative override absence.} There is no administrative role, system operator, or security principal with the capability to modify $\alpha(n)$. This is a deliberate architectural omission. In conventional systems, administrative compromise leads directly to immutability violation because the administrative role holds override authority. In RSP, there is no override authority for $\alpha(n)$ by design.
    \item \textbf{Source-equivalence.} The Fact Layer write protocol treats all modification requests identically regardless of source. A request from the Purpose Layer, the Bounds Engine, an authenticated administrator, or an external adversarial actor receives the same response: protocol-level rejection.
    \item \textbf{Atomic warrant-pointer creation.} The parent pointer and its Purpose Layer warrant are created as a single atomic Fact Layer write. There is no temporal window in which $\alpha(n)$ is established without a matching warrant, or in which a warrant exists without a corresponding parent pointer.
\end{itemize}

\medskip
\noindent\textbf{The sealed memory envelope.} The immutability guarantee is reinforced by the sealed memory envelope mechanism:

\begin{enumerate}
    \item \textbf{Encrypted at rest.} $\alpha(n)$ is encrypted in the node's persistent memory envelope using a Fact Layer-only symmetric key. Even a compromised Bounds Engine cannot read $\alpha(n)$ because it never has access to the decryption key.
    \item \textbf{HMAC-signed.} The envelope carries an HMAC-SHA-256 signature computed over its contents using a Fact Layer-only key. Any modification to the ciphertext invalidates the signature.
    \item \textbf{Decrypted only for reads.} The Fact Layer decrypts the envelope only when reading $\alpha(n)$ to service a chain-traversal or audit request. The plaintext is never exposed to the Bounds Engine or the child node's own code.
    \item \textbf{Re-sealed after reads.} After each authorized read, the Fact Layer re-encrypts and re-signs the envelope, preventing replay attacks based on stale plaintext.
\end{enumerate}

The Fact Layer operates as a standalone, isolated process with its own memory space, inaccessible to the child node's address space. The combination of the protocol-level write constraint and the sealed envelope ensures that $\alpha(n)$ is modification-resistant in all deployment contexts.

\medskip
\noindent\textbf{Why access control is insufficient.} The distinction between access-control-based immutability and protocol-level immutability is the distinction between a promise and a guarantee. In an access-control model, immutability is enforced by restricting which principals may issue modification operations; if an attacker acquires sufficient privileges through credential theft, insider compromise, or software vulnerability, the barrier is circumvented. The immutability was a policy, enforceable until it was not.

In the Fact Layer write protocol, the immutability of $\alpha(n)$ is a property of the protocol itself. There is no operation defined for modifying $\alpha(n)$ after provisioning, so no amount of privilege elevation, credential compromise, or software vulnerability can produce a valid modification operation. The protocol does not enforce immutability — it makes immutability the only valid state transition. The distinction is architectural, not a matter of degree.

This is the precise failure mode the Fact Layer write protocol eliminates: retroactive manipulation of chain topology for the purpose of accountability laundering. Without protocol-level immutability, an adversarial actor could re-parent a node to a favorable parent after the fact, obscuring the actual delegation chain. With the Fact Layer write protocol, this manipulation is structurally impossible because no valid operation exists to produce it.

% ------------------------------------------------------------------------- %
\subsection{Spawn Authorization: Purpose Layer Role}
\label{sec:purpose-layer-authorization}

The Purpose Layer plays two structurally necessary roles in the parent pointer lifecycle: it originates the spawn authorization policy that defines conditions for creating a new parent pointer, and it issues the spawn warrant that is the prerequisite for the Fact Layer provisioning write. These roles are inseparable from the parent pointer's immutability: if the Purpose Layer could issue warrants after the fact, or if a machine actor could create a parent pointer without a Purpose Layer warrant, the chain's integrity guarantees would be void.

\begin{dfn}[Purpose Layer Spawn Authorization]
\label{dfn:purpose-authorization}
A parent node $n_i \in \mathcal{N}$ that intends to spawn a child $n_{i+1}$ must submit a spawn request to the Purpose Layer declaring: (P1) the proposed child operating scope $R(n_{i+1})$, (P2) the specific operational condition necessitating spawning, and (P3) a demonstration that the condition cannot be resolved within $R(n_i)$. The Purpose Layer validates three mandatory conditions before issuing a spawn warrant:

\begin{enumerate}
    \item[(P1)] \textbf{Scope refinement.} $R(n_{i+1})$ must be a strict subset of $R(n_i)$:
    \begin{equation}
    R(n_{i+1}) \subset R(n_i).
    \label{eq:scope-refinement}
    \end{equation}
    Lateral scope stability — where $R(n_{i+1}) = R(n_i)$ — is not authorized. The chain grows through scope \emph{narrowing}, not sideways expansion. Equality at any link constitutes a drift condition and invalidates the warrant.

    \item[(P2)] \textbf{Operational necessity.} The parent must declare, in the Fact Layer authorization ledger, the precise condition requiring child spawning. The declaration must identify why the condition cannot be resolved within $R(n_i)$ and why the child is the minimum necessary decomposition of the current task.

    \item[(P3)] \textbf{Minimum scope proportionality.} The child scope must be the narrowest scope sufficient to address the declared condition. The Purpose Layer does not authorize child scopes broader than operationally necessary.
\end{enumerate}
\end{dfn}

When all three conditions are satisfied, the Purpose Layer issues a spawn warrant $w_i$ — a cryptographically signed record written atomically to the Fact Layer authorization ledger:

\begin{equation}
w_i = \bigl\langle
\text{parent}=n_i,\;
\text{child}=n_{i+1},\;
R(n_{i+1}),\;
\text{justification},\;
\text{timestamp},\;
\sigma_{\text{PL}}
\bigr\rangle .
\label{eq:spawn-warrant}
\end{equation}

The warrant is the sole authorized precondition for the Fact Layer's $\mathrm{WriteParentPointer}$ operation. No machine actor may initiate a spawn event without a matching warrant in the Fact Layer ledger — the Fact Layer pre-commit hook rejects any provisioning request that lacks a valid warrant before the operation reaches the ledger write stage. This makes the Purpose Layer's exclusive terminus role a structural guarantee, not a configuration option.

The Purpose Layer also enforces two structural constraints at spawn authorization:

\begin{enumerate}
    \item[(C1)] \textbf{Cycle prevention.} The proposed parent $n_i$ must not be a descendant of the proposed child $n_{i+1}$ in the current chain. This is verified by checking $n_{i+1} \notin \Chain{n_i}$. If the proposed child already appears in the parent's ancestry chain, the spawn is denied because it would introduce a cycle in $\alpha$, violating the acyclicity of the delegation graph.
    \item[(C2)] \textbf{Depth bound compliance.} The depth of the proposed child must satisfy $\mathrm{depth}(n_{i+1}) \leq D_{\max}$, where $D_{\max}$ is the architectural maximum depth parameter recorded in the Fact Layer authorization ledger at system initialization. If the spawn would cause the depth bound to be exceeded, the spawn is refused and the node is placed in \texttt{ESCALATING} state with a mandatory Bailout Protocol trigger.
\end{enumerate}

The Purpose Layer's exclusive authority over spawn authorization means that no machine actor in the chain can autonomously extend the chain. Every link in the delegation chain is Purpose Layer-authorized before it is Fact Layer-recorded. This two-layer gate — Purpose Layer warrant first, Fact Layer recording second — is the fundamental control structure that makes the RSP's accountability guarantees operative.

% ------------------------------------------------------------------------- %
\subsection{Chain Traversal Algorithm}
\label{sec:chain-traversal}

The chain traversal algorithm provides the runtime mechanism for computing $\Chain{a}$ from any node $a$, and the chain verification algorithm provides the runtime procedure for confirming the structural validity of the chain to a human principal.

\begin{algorithm}
\caption{Chain-Traverse($n$) — Build the delegation chain from node $n$ to the human anchor}
\label{alg:chain-traverse}
\begin{algorithmic}[1]
\REQUIRE $n$ is a provisioned RSP node with $\alpha$ defined
\ENSURE $\Chain{n}$ — the ordered list of all ancestor nodes including the human anchor
\STATE $chain \leftarrow \langle \rangle$ \COMMENT{Initialize empty ordered list}
\STATE $current \leftarrow n$ \COMMENT{Start at the query node}
\WHILE{$\mathrm{true}$}
    \STATE $\mathrm{append}(chain,\ current)$ \COMMENT{Add current node to chain}
    \IF{$\alpha(current) = \bot$}
        \STATE $\textbf{break}$ \COMMENT{Root human anchor reached — chain complete}
    \ENDIF
    \STATE $current \leftarrow \alpha(current)$ \COMMENT{Move to parent}
\ENDWHILE
\RETURN $chain$
\end{algorithmic}
\end{algorithm}

\begin{algorithm}
\caption{Chain-Verify($n$) — Verify the authorization chain from node $n$ to human anchor}
\label{alg:chain-verify}
\begin{algorithmic}[1]
\REQUIRE $n$ is a provisioned RSP node
\ENSURE $\mathrm{Boolean}$ indicating whether the chain is a structurally valid RSP chain
\STATE $chain \leftarrow \mathrm{Chain\text{-}Traverse}(n)$
\STATE $m \leftarrow \mathrm{len}(chain)$

\STATE \COMMENT{\textit{Step V1: Verify chain terminates at human anchor}}
\IF{$\alpha(chain[m-1]) \neq \bot$}
    \STATE \RETURN $\text{false}$ \COMMENT{Root is not $\bot$ — malformed chain}
\ENDIF
\IF{$\neg h(chain[m-1])$}
    \STATE \RETURN $\text{false}$ \COMMENT{Root is not a designated human principal}
\ENDIF

\STATE \COMMENT{\textit{Step V2: Verify each spawn event has a Purpose Layer warrant}}
\FOR{$i \leftarrow 0$ \TO $m-2$}
    \STATE $child \leftarrow chain[i]$
    \STATE $parent \leftarrow chain[i+1]$
    \IF{$\neg \text{warrant\_exists}(child,\ parent)$}
        \STATE \RETURN $\text{false}$ \COMMENT{Missing Purpose Layer warrant for spawn}
    \ENDIF
\ENDFOR

\STATE \COMMENT{\textit{Step V3: Verify scope refinement at each link}}
\FOR{$i \leftarrow 0$ \TO $m-2$}
    \STATE $child \leftarrow chain[i]$
    \STATE $parent \leftarrow chain[i+1]$
    \IF{$\neg (R(child) \subset R(parent))$}
        \STATE \RETURN $\text{false}$ \COMMENT{Scope refinement constraint violated}
    \ENDIF
\ENDFOR

\STATE \COMMENT{\textit{Step V4: Verify depth bound compliance at each node}}
\FOR{$i \leftarrow 0$ \TO $m-1$}
    \IF{$\mathrm{depth}(chain[i]) > D_{\max}$}
        \STATE \RETURN $\text{false}$ \COMMENT{Depth bound exceeded at node}
    \ENDIF
\ENDFOR

\RETURN $\text{true}$
\end{algorithmic}
\end{algorithm}

Chain-Verify returns \texttt{true} if and only if all four verification conditions hold simultaneously. V1 confirms the chain terminates at a human anchor; V2 confirms every spawn was authorized by the Purpose Layer; V3 confirms scope refinement was respected at every link; V4 confirms the depth bound was respected throughout. A chain that passes Chain-Verify is, by the RSP definitions, a fully authorized delegation chain terminating at a human principal.

The algorithm is deterministic and produces the same result regardless of which node initiates verification. There is no self-serving interpretation of authorization that a drifted or orphaned node could produce, because Chain-Verify inspects the complete chain topology and warrants in the Fact Layer ledger, not the node's own claims about its authorization.

\medskip
\noindent\textbf{Correctness arguments.}

\textit{Termination.} The algorithm terminates because RSP chains are depth-bounded. Each iteration of the while-loop moves $current$ to $\alpha(current)$, which is strictly closer to the root. Since the chain has finite depth at most $D_{\max} + 1$, the loop executes at most $D_{\max} + 1$ iterations before $\alpha(current) = \bot$ is reached. The loop therefore terminates in finite time regardless of the operational state of any node in the chain.

\textit{Correctness of the returned set.} By Definition~\ref{dfn:chain-operator}, $\Chain{a}$ is the set containing $a$, $\alpha(a)$, $\alpha(\alpha(a))$, and so on, until the root. The algorithm produces exactly this sequence: it starts at $a$, appends each node, moves to its parent via $\alpha$, and stops when $\alpha(current) = \bot$ is reached. The output is identical to $\Chain{a}$ by construction.

\textit{Uniqueness of result.} The algorithm produces a unique result for any input node because $\alpha$ is a total function (defined for every provisioned node) and injective on the chain direction (every node has exactly one parent, and the root has no parent). There is no branching, no alternative path, and no non-deterministic choice at any step.

\textit{Immutability guarantee.} The algorithm reads $\alpha$ values but never writes them. The traversal is read-only and does not interact with the Fact Layer write protocol. The algorithm can be executed at any runtime point without affecting the immutability of the chain it traverses, and without creating any modification opportunities for adversarial actors.

\medskip
\noindent\textbf{Constant-time escalation bound.} The depth bound $D_{\max}$ provides a constant-time upper bound on chain traversal cost. The worst-case number of pointer dereferences required to reach the human anchor from any node is $D_{\max} + 1$, independent of the total number of nodes in the system. A system with ten thousand spawned nodes and $D_{\max} = 5$ requires at most six pointer dereferences to reach a human — the same as a system with ten nodes. Mutable parent pointers would allow the chain to be extended retroactively, making the effective depth potentially unbounded and the worst-case escalation time unbounded as well. Immutability is what makes the depth bound a structural guarantee, not merely a policy.

In the Bailout Protocol context, chain traversal is used to identify the correct Human Accountability Anchor when an agent at any depth encounters an unresolvable exception. The upward propagation of the exception signal follows the same parent-pointer structure as chain traversal, but uses the asynchronous trigger pathway rather than the synchronous verification pathway. Chain traversal ensures that this walk never diverges, never shortcuts, and never fails to reach a human.

\medskip
\noindent\textbf{Computational cost.} Chain-Traverse runs in $O(d)$ time where $d$ is the chain depth (at most $D_{\max} + 1$). Chain-Verify runs in $O(m)$ for the traversal plus $O(m)$ for each of the three verification loops, yielding $O(m)$ total where $m$ is chain length, also bounded by $D_{\max} + 1$. Both algorithms are linear in chain length and do not depend on the total number of nodes in the system.

% ------------------------------------------------------------------------- %
\subsection{Summary: Immutability as Architectural Constraint}
\label{sec:immutability-summary}

The immutable parent pointer is the load-bearing constraint of the Recursive Spawning Protocol. Without it, the chain is mutable — subject to retroactive modification by actors with sufficient privilege — and the RSP's correctness properties (drift prevention, chain integrity, termination) collapse to policy conventions rather than structural guarantees.

The immutability is architectural — not a stronger access control policy — because it is enforced by the Fact Layer write protocol, which defines no valid modification operation for $\alpha(n)$ after provisioning. This is distinct from access-control-based immutability in four ways: there is no override path for any actor; the immutability holds under all system failure conditions including those that disable other enforcement mechanisms; all sources receive equivalent enforcement treatment regardless of privilege level; and the parent pointer and its Purpose Layer warrant are created atomically with no temporal window for inconsistency.

The chain operator $\Chain{a}$ is uniquely determined by the immutable parent pointers in the chain. The base case of the root human anchor ($\alpha(h) = \bot$) terminates the chain and is the only valid terminus. The chain traversal algorithm computes this operator correctly and terminates. Chain-Verify provides the runtime verification procedure for chain structural validity. Together, these mechanisms make the delegation chain an immutable, verifiable, and auditable runtime artifact — the foundation on which all other RSP correctness properties rest.
